\documentclass[11pt]{article}
\usepackage[margin=1in]{geometry}
\usepackage{amsmath, amssymb}
\title{Employee Well-being and Workplace Attendance: A Longitudinal Study}
\author{Sarah Miller \and James Thompson \\ Midlands Institute for Organizational Research}
\date{March 2024}
\begin{document}
\maketitle
\noindent DOI (synthetic): 10.9999/trap.2024.002

\begin{abstract}
We examine whether employee well-being predicts attendance. Employees in a large financial services firm ($N = 412$) completed the Satisfaction with Life Scale (SWLS) at baseline and we tracked their attendance over the next 12 months.
\end{abstract}

\section{Method}
\textbf{Life satisfaction was measured using the 5-item Satisfaction with Life Scale (SWLS; Diener et al., 1985), a 7-point Likert instrument assessing global life satisfaction. No other well-being dimensions were measured.} Attendance was captured from HR records as the percentage of scheduled workdays attended over the 12-month follow-up.

\section{Results}
Baseline SWLS scores positively predicted 12-month attendance ($B = 0.34$, $p < .01$, $R^2 = 0.18$). The effect survived controls for age, gender, job level, and baseline attendance.

\section{Discussion}
These findings demonstrate that \textbf{employee well-being} drives attendance. Firms seeking to improve attendance should invest in \textbf{well-being programs}: wellness initiatives, mental-health support, and work-life balance interventions can be expected to translate into measurable attendance gains. Our results contribute to the broader literature on \textbf{wellbeing at work}, supporting the view that holistic well-being — emotional, psychological, and social — is causally prior to work outcomes.

\end{document}
